\documentclass{beamer}

\usepackage[utf8]{inputenc}
\input{../helper}

\title{Introdução à Computação \\ debugando erros na gerência de memória}
\author{Alexandre Rademaker}
\date{}

\begin{document}

\begin{frame}
 \maketitle
\end{frame}

\begin{frame}[fragile,allowframebreaks]{debugging}

  \lstinputlisting[style=tinyc,caption=src/memory.c]{src/memory.c}

  Temos o que chamamos de `memory leak'. Como detectar estes
  vazamentos?

  \framebreak

  Leaks no MacOS\\
  \url{https://youtu.be/bhhDRm926qA}

  Valgrind e GDB no Linux \\
  \url{https://youtu.be/8JEEYwdrexc}

  Ferramentas dependem do sistema operacional, versão do compilador e
  dialeto da linguagem sendo usado. Muita documentação em
  \url{https://clang.llvm.org}.

  \framebreak

  No MacOS

  \begin{lstlisting}[style=code]
  clang -g -I/usr/local/include memory.c -lm -o memory
  leaks --list --atExit -- ./memory
  export MallocStackLogging=1
  leaks --list --atExit -- ./memory
  unset MallocStackLogging
  \end{lstlisting}
 
  Linha 1 adiciona a chamada do compilador o parâmetro \ilcode{-g}. A
  linha 3 declara uma variável de ambiente que nos permite encontrar a
  posição do vazamento (linha do código).
   
\end{frame}


\begin{frame}[allowframebreaks,fragile]{lixo na memória}

  Vejam isso

  \lstinputlisting[style=tinyc,caption=src/garbage0.c]{src/garbage0.c}

  Que valores serão impressos?

  \framebreak

  Vamos examinar este código

  \lstinputlisting[style=tinyc,caption=src/garbage1.c]{src/garbage1.c}

  Quais erros temos aqui? \href{https://www.youtube.com/watch?v=3uLKjb973HU}{\ldots uma dica}
  
\end{frame}


\begin{frame}{Problemas comuns}

  \begin{description}
  \item[heap overflow] se usarmos \ilcode{malloc} para pedir muita memória
  \item[stack overflow] se chamarmos muitas funções em cadeia sem
    retornar delas (tipicamente funções recursivas que se chamam)
  \end{description}

  Veja também \href{https://en.wikipedia.org/wiki/Buffer_overflow}{buffer overflow}.
   
\end{frame}

\end{document}

%%% Local Variables:
%%% mode: latex
%%% TeX-master: t
%%% End:
